\documentclass[10pt]{article}
\usepackage[left=0.8in,right=0.8in,top=0.8in,bottom=0.8in]{geometry}
\usepackage{float}
\input{../notation}
\usepackage[ruled,lined]{algorithm2e}
\usepackage{titling}
\setlength{\droptitle}{-2em}
\posttitle{\par\end{center}\vspace{-1em}}
\setlength{\parskip}{0.3em}
\setlength{\parindent}{0em}

\title{Primer: Temporal-Difference CCP (TD-CCP)\\[4pt]
\large Adusumilli and Eckardt (2025)}
\author{econirl package documentation}
\date{}

\begin{document}
\maketitle

%% ============================================================
\section{Context and Problem}
%% ============================================================

CCP estimation (Hotz and Miller 1993) avoids the inner Bellman loop of NFXP by inverting observed choice probabilities into value differences.
The inversion requires computing the expected future value $\E[V(s') \mid s,a]$, which in turn requires the transition density $p(s' \mid s,a)$.
When the state space is continuous or high-dimensional, estimating this density is the binding constraint.
Kernel methods suffer from the curse of dimensionality, and discretization introduces approximation bias that grows with the number of state variables.

Adusumilli and Eckardt (2025) eliminate the transition density entirely.
Their key insight is that the CCP value function components $h(a,x)$ and $g(a,x)$ satisfy temporal-difference fixed-point equations that depend only on sample expectations over observed transitions $(x_t, a_t, x_{t+1})$, not on the density $p(x' \mid x,a)$.
This is the same idea that makes TD learning in reinforcement learning model-free: replace $\E_{p}[V(s')]$ with the sample next state $V(s'_{t+1})$.

The estimator comes in two flavors.
The linear semi-gradient method approximates $h$ and $g$ with polynomial basis functions and solves a single linear system, running in sub-second time.
The neural approximate value iteration (AVI) method uses neural networks for $h$ and $g$, iterating Bellman-like regression targets until convergence.
Both produce locally robust standard errors at the parametric $\sqrt{n}$ rate, even though the first-stage $h$ and $g$ estimates converge at slower nonparametric rates.

TD-CCP is the only estimator in econirl that requires neither transition densities nor an inner Bellman loop.
NFXP and CCP require known transitions.
NNES and SEES require transitions for the Bellman target.
AIRL avoids transitions but does not produce valid standard errors.
TD-CCP provides all three: no transitions, scalability, and valid inference.

The name ``temporal-difference CCP'' reflects the connection to TD learning in reinforcement learning.
Just as TD($0$) updates value estimates using the observed next state rather than the full expectation over the transition kernel, TD-CCP updates value function components $h_k$ and $g$ using observed next transitions $(x_{t+1}, a_{t+1})$ rather than the density $p(x' \mid x,a)$.
The semi-gradient fixed-point equation~\eqref{eq:semigradient} is the structural econometrics analog of the Tsitsiklis and Van Roy (1997) linear TD method with function approximation.

%% ============================================================
\section{Model and Notation}
%% ============================================================

\begin{table}[h!]
\centering\small
\begin{tabular}{ll}
\toprule
Symbol & Meaning \\
\midrule
$x \in \mathcal{X}$ & State (continuous or discrete) \\
$a \in \mathcal{A}$ & Action (discrete, $|\mathcal{A}|$ choices) \\
$\beta \in [0,1)$, $\sigma > 0$ & Discount factor, logit scale \\
$\phi_k(a,x) \in \R$ & Feature $k$ for action $a$ at state $x$ \\
$\theta \in \R^K$ & Structural parameters \\
\midrule
$P(a \mid x)$ & Conditional choice probability (CCP) \\
$e(a,x) = \gamma_E - \ln P(a \mid x)$ & Emax correction (Euler constant $\gamma_E \approx 0.577$) \\
$h_k(a,x)$ & Expected discounted feature $k$ under action $a$ \\
$g(a,x)$ & Expected discounted emax under action $a$ \\
$z_k(a,x), z_g(a,x)$ & Pseudo-outcome for $h_k$, $g$ \\
\midrule
$\Psi(a,x) \in \R^{K_\Psi}$ & Basis functions for semi-gradient \\
$\omega_k \in \R^{K_\Psi}$ & Basis coefficients for $h_k$ \\
$\lambda_k(a,x)$ & Backward value function (for robust SEs) \\
\bottomrule
\end{tabular}
\end{table}
\FloatBarrier

The CCP decomposition of Hotz and Miller (1993) writes the choice-specific value as
\begin{equation}
\label{eq:ccp_decomp}
Q(a,x;\theta) \;=\; \sum_{k=1}^K \theta_k\,h_k(a,x) \;+\; \sigma\,g(a,x),
\end{equation}
where $h_k(a,x) = \E\bigl[\sum_{t=0}^\infty \beta^t \phi_k(a_t, x_t) \mid a_0 = a, x_0 = x\bigr]$ is the expected discounted feature $k$ and $g(a,x) = \E\bigl[\sum_{t=1}^\infty \beta^t e(a_t, x_t) \mid a_0 = a, x_0 = x\bigr]$ is the expected discounted emax.
The pseudo-log-likelihood is
\begin{equation}
\label{eq:pseudo_ll}
Q_n(\theta) \;=\; \frac{1}{n}\sum_{i=1}^n \ln \pi\!\bigl(a_i \mid x_i;\, \theta,\, h,\, g\bigr), \qquad \pi(a \mid x) = \frac{\exp\bigl(Q(a,x;\theta)/\sigma\bigr)}{\sum_{a'}\exp\bigl(Q(a',x;\theta)/\sigma\bigr)}.
\end{equation}
The key objects $h_k$ and $g$ satisfy temporal-difference fixed-point equations that do not involve the transition density.

%% ============================================================
\section{Identification and Theory}
%% ============================================================

\begin{proposition}[TD fixed point for $h$]
The component $h_k(a,x)$ satisfies the fixed-point equation $h_k(a,x) = z_k(a,x) + \beta\,\E[h_k(a',x') \mid a,x]$, where $z_k(a,x) = \phi_k(a,x) + \beta\sum_{a'} P(a' \mid x)\bigl[\phi_k(a',x) - \phi_k(a,x)\bigr]$ is the pseudo-outcome and $(a',x')$ is the next observed state-action pair.
The expectation is over the data-generating process, not over a model transition density.
\end{proposition}

The semi-gradient method approximates $h_k(a,x) \approx \Psi(a,x)^\top \omega_k$ and solves
\begin{equation}
\label{eq:semigradient}
\hat\omega_k \;=\; \Bigl[\frac{1}{n}\sum_{i=1}^n \Psi_i\bigl(\Psi_i - \beta\,\Psi_i'\bigr)^\top\Bigr]^{-1} \frac{1}{n}\sum_{i=1}^n \Psi_i\, z_{k,i},
\end{equation}
where $\Psi_i = \Psi(a_i, x_i)$ and $\Psi_i' = \Psi(a_i', x_i')$ are basis evaluations at current and next transitions.
This is a single matrix inversion, shared across all $K$ components.

\begin{proposition}[Locally robust inference]
Define the backward value function $\lambda_k(a,x)$ as the solution to the adjoint equation $(I - \beta F_\pi^\top)\lambda_k = \partial Q / \partial h_k$.
The locally robust moment condition is
\begin{equation}
\label{eq:robust_moment}
\zeta_i(\theta) \;=\; m_i(\theta) \;+\; \sum_k \theta_k\,\lambda_k(a_i', x_i')\,\delta_{h,k,i} \;+\; \sigma\,\lambda_g(a_i', x_i')\,\delta_{g,i},
\end{equation}
where $m_i$ is the standard score, $\delta_{h,k,i} = z_{k,i} + \beta\,\hat h_k(a_i', x_i') - \hat h_k(a_i, x_i)$ is the TD error for $h_k$, and $\delta_{g,i}$ is defined analogously for $g$.
Under regularity conditions, $\hat\theta$ is asymptotically normal at rate $\sqrt{n}$ with sandwich variance $V = (G^\top \Omega^{-1} G)^{-1}$, even when $\hat h$ and $\hat g$ converge at slower nonparametric rates.
\end{proposition}

The key property is Neyman orthogonality: the moment condition $\zeta_i$ has zero derivative with respect to the nuisance functions $h$ and $g$ at the true values.
This means first-order errors in $\hat h$ and $\hat g$ drop out of the influence function for $\hat\theta$, so the structural parameters inherit the parametric $\sqrt{n}$ rate even when the nonparametric first stage converges at a slower rate such as $n^{-1/3}$ or $n^{-2/5}$.
Without the correction terms in~\eqref{eq:robust_moment}, naive standard errors would understate uncertainty by ignoring the estimation error in $h$ and $g$.

Cross-fitting (splitting data into two folds, estimating $h,g$ on one fold and $\theta$ on the other) eliminates the Donsker condition that would otherwise restrict the complexity of the first-stage estimator.
On tabular problems where the state space is finite, the standard approach without cross-fitting also produces valid SEs because the CCP estimates converge at the parametric rate.
The paper recommends cross-fitting for continuous-state applications where nonparametric rates apply.

The identification conditions for $\theta$ are the same as for NFXP: the feature matrix must have full column rank, and the discount factor $\beta$ must be fixed a priori.
The TD-CCP approach does not weaken or strengthen the identification requirements relative to standard CCP.
What it changes is the computational requirements: the transition density is no longer needed for estimation, only for counterfactual simulation after estimation.

%% ============================================================
\section{Estimation Algorithm}
%% ============================================================

\begin{algorithm}[H]
\DontPrintSemicolon
\caption{TD-CCP with Semi-Gradient (Adusumilli and Eckardt 2025)}
\label{alg:tdccp}
\KwIn{Panel $\{(x_i, a_i, x_i', a_i')\}_{i=1}^n$, features $\phi$, basis $\Psi$, discount $\beta$}
\textbf{Cross-fitting:} split data into folds $\mathcal{I}_1, \mathcal{I}_2$\tcp*{\texttt{\_estimate\_with\_cross\_fitting()}}
\For{fold $\ell = 1, 2$}{
  \textbf{Step 1: Estimate CCPs}\tcp*{\texttt{\_estimate\_ccps()}}
  $\hat P(a \mid x) \leftarrow$ frequency counts (or logit on polynomial features)\;
  $e_i \leftarrow \gamma_E - \ln \hat P(a_i \mid x_i)$ for each observation\;
  \BlankLine
  \textbf{Step 2: Build pseudo-outcomes}\;
  $z_{k,i} \leftarrow \phi_k(a_i,x_i) + \beta\sum_{a'} \hat P(a' \mid x_i)[\phi_k(a',x_i) - \phi_k(a_i,x_i)]$\;
  $z_{g,i} \leftarrow e_i + \beta\sum_{a'} \hat P(a' \mid x_i)[e(a',x_i) - e(a_i,x_i)]$\;
  \BlankLine
  \textbf{Step 3: Solve semi-gradient}\tcp*{\texttt{\_semigradient\_solve()}}
  $A \leftarrow \frac{1}{|\mathcal{I}_\ell|}\sum_{i \in \mathcal{I}_\ell} \Psi_i(\Psi_i - \beta\Psi_i')^\top$\;
  \For{$k = 1, \ldots, K$ and for $g$}{
    $\hat\omega_k \leftarrow A^{-1}\bigl[\frac{1}{|\mathcal{I}_\ell|}\sum_{i \in \mathcal{I}_\ell} \Psi_i\,z_{k,i}\bigr]$\tcp*{Eq.~\ref{eq:semigradient}}
  }
  \BlankLine
  \textbf{Step 4: Partial MLE on other fold}\tcp*{\texttt{\_partial\_mle()}}
  $\hat\theta^{(\ell)} \leftarrow \argmax_\theta \frac{1}{|\mathcal{I}_{-\ell}|}\sum_{i \in \mathcal{I}_{-\ell}} \ln \pi(a_i \mid x_i;\,\theta,\,\hat h,\,\hat g)$\;
  \BlankLine
  \textbf{Step 5: Backward value function}\tcp*{\texttt{\_compute\_backward\_value()}}
  Solve $(I - \beta\hat F_\pi^\top)\hat\lambda_k = \partial Q/\partial h_k$ for each $k$\;
  \BlankLine
  \textbf{Step 6: Robust SEs}\tcp*{\texttt{\_compute\_robust\_se()}}
  Compute $\hat\zeta_i$ via~\eqref{eq:robust_moment}\;
  $\hat V^{(\ell)} \leftarrow (\hat G^\top \hat\Omega^{-1} \hat G)^{-1}$\;
}
$\hat\theta \leftarrow \tfrac{1}{2}(\hat\theta^{(1)} + \hat\theta^{(2)})$, \quad $\hat V \leftarrow \tfrac{1}{2}(\hat V^{(1)} + \hat V^{(2)})$\;
\KwOut{$\hat\theta$, $\hat V$ for robust SEs, $\hat h$, $\hat g$ for diagnostics}
\end{algorithm}

Code annotations reference methods in \texttt{src/econirl/estimation/td\_ccp.py}.
The matrix $A$ in Step~3 is shared across all components, so the cost is one matrix inversion plus $K+1$ back-solves.
Default configuration: $K_\Psi = 8$ polynomial basis, 2-fold cross-fitting, robust SEs enabled.
For neural AVI, replace Step~3 with iterative regression networks (\texttt{\_neural\_avi\_solve()}).

%% ============================================================
\section{Simulation}
%% ============================================================

% Auto-generated by tdccp_run.py. Do not edit by hand.
% Known-truth DGP cells: canonical_low_action, canonical_high_action

\section{Known-Truth Validation Results}
TD-CCP is transition-free for structural parameter estimation: the estimator uses observed state-action-next-state tuples to estimate the recursive terms \(h(a,x)\) and \(g(a,x)\). Known transitions enter only after estimation, when the recovered reward is evaluated for policies, values, and counterfactuals.

The paper-faithful hard flexible DGP is the certified release path: it uses two-dimensional encoded state coordinates, logit first-stage CCPs, Algorithm 2 cross-fitting, locally robust standard errors, and a finite linear structural reward parameter. The low-dimensional action-dependent cell is retained as a sanity check. The high-dimensional encoded-state cell and raw neural reward cell are diagnostic stress tests, not release certification.

\subsection{Validation Design Context}
\begin{itemize}
\item \textbf{low-dimensional action-dependent DGP.} A compact finite-state DDC benchmark with action-specific rewards, known transitions, known reward, value, policy, and Q-function truth, and Type A/B/C counterfactual oracles.
\item \textbf{high-dimensional action-dependent DGP.} A harder encoded-state and richer reward-feature stress cell using the same known-truth reward, value, policy, Q-function, and Type A/B/C counterfactual objects to test scaling beyond the compact tabular case.
\item \textbf{encoded-state locally robust DGP.} A TD-CCP hard cell with two-dimensional encoded state coordinates, stochastic transitions, logit first-stage CCPs, Algorithm 2 cross-fitting, locally robust standard errors, and a finite linear structural reward with action 0 normalized to zero.
\item \textbf{raw-neural flexible diagnostic DGP.} A raw-neural flexible diagnostic with a frozen neural reward matrix and no finite true theta, so it is not a finite-parameter recovery claim.
\end{itemize}

\subsection{Validation Cells}
\begin{table}[H]
\centering\scriptsize
\caption{TD-CCP known-truth cells.}
\begin{tabular}{@{}llrrrrrr@{}}
\toprule
Cell & Role & States & State dim. & Reward params & Iter. & Gates pass & Gates fail \\
\midrule
low-dimensional action-dependent DGP & low-dimensional check & 21 & 2 & 4 & 14 & 10 & 0 \\
high-dimensional action-dependent DGP & high-dimensional diagnostic & 81 & 16 & 32 & 23 & 0 & 10 \\
\bottomrule
\end{tabular}
\end{table}

\subsection{Estimator Settings}
\begin{table}[H]
\centering\scriptsize
\caption{TD-CCP settings used by the shared harness.}
\begin{tabular}{lrrrrr}
\toprule
Cell & Basis & CCP smoothing & L2 penalty & Cross-fit & Transitions for theta \\
\midrule
low-dimensional action-dependent DGP & encoded-3 & 0.01 & 0 & false & no \\
high-dimensional action-dependent DGP & encoded-3 & 0.10 & 1000 & false & no \\
\bottomrule
\end{tabular}
\end{table}

\subsection{Recovery Metrics}
\begin{table}[H]
\centering\scriptsize
\caption{Core structural recovery metrics.}
\begin{tabular}{lrrrrrr}
\toprule
Cell & Param. cos. & Param. rel. RMSE & Reward RMSE & Policy TV & Value RMSE & Q RMSE \\
\midrule
low-dimensional action-dependent DGP & 0.998 & 0.078 & 0.012 & 0.007 & 0.025 & 0.029 \\
high-dimensional action-dependent DGP & 0.939 & 0.346 & 0.164 & 0.069 & 0.905 & 0.742 \\
\bottomrule
\end{tabular}
\end{table}

\subsection{Gate Audit}
\begin{table}[H]
\centering\scriptsize
\caption{Hard-gate audit.}
\begin{tabular}{llrrr}
\toprule
Cell & Gate & Value & Threshold & Pass \\
\midrule
low-dimensional action-dependent DGP & converged & true & is true & true \\
low-dimensional action-dependent DGP & parameter\_cosine & 0.998 & >= 0.990 & true \\
low-dimensional action-dependent DGP & parameter\_relative\_rmse & 0.078 & <= 0.150 & true \\
low-dimensional action-dependent DGP & reward\_rmse & 0.012 & <= 0.030 & true \\
low-dimensional action-dependent DGP & policy\_tv & 0.007 & <= 0.020 & true \\
low-dimensional action-dependent DGP & value\_rmse & 0.025 & <= 0.100 & true \\
low-dimensional action-dependent DGP & q\_rmse & 0.029 & <= 0.100 & true \\
low-dimensional action-dependent DGP & type\_a\_regret & 0.000 & <= 0.010 & true \\
low-dimensional action-dependent DGP & type\_b\_regret & 0.000 & <= 0.010 & true \\
low-dimensional action-dependent DGP & type\_c\_regret & 0.000 & <= 0.010 & true \\
high-dimensional action-dependent DGP & converged & false & is true & false \\
high-dimensional action-dependent DGP & parameter\_cosine & 0.939 & >= 0.990 & false \\
high-dimensional action-dependent DGP & parameter\_relative\_rmse & 0.346 & <= 0.150 & false \\
high-dimensional action-dependent DGP & reward\_rmse & 0.164 & <= 0.030 & false \\
high-dimensional action-dependent DGP & policy\_tv & 0.069 & <= 0.020 & false \\
high-dimensional action-dependent DGP & value\_rmse & 0.905 & <= 0.100 & false \\
high-dimensional action-dependent DGP & q\_rmse & 0.742 & <= 0.100 & false \\
high-dimensional action-dependent DGP & type\_a\_regret & 0.318 & <= 0.010 & false \\
high-dimensional action-dependent DGP & type\_b\_regret & 0.593 & <= 0.010 & false \\
high-dimensional action-dependent DGP & type\_c\_regret & 0.015 & <= 0.010 & false \\
\bottomrule
\end{tabular}
\end{table}

\subsection{Encoded-State Locally Robust Showcase}
The certified hard DGP uses stochastic transitions and two-dimensional encoded state coordinates. The reward is linear in an intercept and the two state coordinates, interacted with nonbaseline actions; action 0 is the utility normalization. The estimator uses logit CCPs, an encoded semi-gradient basis, Algorithm 2 cross-fitting, and locally robust standard errors.

The encoded-state locally robust DGP passes all gates. This is the condition under which the RTD page may describe the finite-theta TD-CCP hard-case showcase as validated.

\begin{table}[H]
\centering\scriptsize
\caption{TD-CCP encoded-state locally robust DGP settings and gate summary.}
\begin{tabular}{llrrrr}
\toprule
Case & Reward / features & Obs. & Params & Gates pass & Gates fail \\
\midrule
encoded-state locally robust DGP & linear / encoded & 120000 & 6 & 14 & 0 \\
\bottomrule
\end{tabular}
\end{table}

\begin{table}[H]
\centering\scriptsize
\caption{TD-CCP encoded-state locally robust DGP fit summary.}
\begin{tabular}{lrrrrr}
\toprule
Case & Converged & Iter. & Log likelihood & Est. time & Obs. \\
\midrule
encoded-state locally robust DGP & true & 121 & -99831.3015 & 10.71 seconds & 120000 \\
\bottomrule
\end{tabular}
\end{table}

\begin{table}[H]
\centering\scriptsize
\caption{TD-CCP locally robust inference diagnostics.}
\begin{tabular}{lrrrrrrr}
\toprule
Case & Cov. unit & Prelim. stat. & Prelim. pgrad & Robust OK & Max. zeta norm & Lambda max abs & Max SE \\
\midrule
encoded-state locally robust DGP & individual & 2/2 & 0.000 & 2/2 & 0.000 & 1.023 & 0.034 \\
\bottomrule
\end{tabular}
\end{table}

\begin{table}[H]
\centering\scriptsize
\caption{Hard flexible DGP recovery metrics. Normalized RMSE divides by the truth RMS.}
\begin{tabular}{lrrrrrrr}
\toprule
Case & Param. cos. & Param. rel. RMSE & Reward nRMSE & Policy TV & Value nRMSE & Q nRMSE \\
\midrule
encoded-state locally robust DGP & 0.999 & 0.059 & 0.025 & 0.005 & 0.001 & 0.001 \\
\bottomrule
\end{tabular}
\end{table}

\begin{table}[H]
\centering\scriptsize
\caption{Hard flexible DGP counterfactual regret checks.}
\begin{tabular}{lrrr}
\toprule
Counterfactual & Policy TV & Value RMSE & Regret \\
\midrule
Type A & 0.005 & 0.002 & 0.002 \\
Type B & 0.005 & 0.002 & 0.002 \\
Type C & 0.007 & 0.003 & 0.003 \\
\bottomrule
\end{tabular}
\end{table}

\begin{table}[H]
\centering\scriptsize
\caption{Hard flexible DGP gate audit.}
\begin{tabular}{lrrr}
\toprule
Gate & Value & Threshold & Pass \\
\midrule
converged & true & is true & true \\
algorithm2\_locally\_robust\_path & true & is true & true \\
finite\_positive\_standard\_errors & true & is true & true \\
zeta\_moment\_norm & 0.000 & <= 0.000 & true \\
covariance\_min\_eigenvalue & 0.000 & >= -0.000 & true \\
parameter\_cosine & 0.999 & >= 0.990 & true \\
parameter\_relative\_rmse & 0.059 & <= 0.150 & true \\
reward\_normalized\_rmse & 0.025 & <= 0.080 & true \\
policy\_tv & 0.005 & <= 0.030 & true \\
value\_normalized\_rmse & 0.001 & <= 0.100 & true \\
q\_normalized\_rmse & 0.001 & <= 0.100 & true \\
type\_a\_regret & 0.002 & <= 0.050 & true \\
type\_b\_regret & 0.002 & <= 0.050 & true \\
type\_c\_regret & 0.003 & <= 0.050 & true \\
\bottomrule
\end{tabular}
\end{table}

\subsection{Repeated-Seed Standard Error Coverage}
The coverage check repeatedly simulates the certified encoded-state DGP and re-estimates TD-CCP with the same Algorithm 2 settings. Reported standard errors use individual-clustered fold covariances. Preliminary plug-in optimizer status is reported separately because the final robust zeta solve is the binding convergence check.

\begin{table}[H]
\centering\scriptsize
\caption{TD-CCP repeated-seed coverage summary.}
\begin{tabular}{lrrrrrr}
\toprule
Case & Reps. & Obs./rep. & Coverage & Mean rel. RMSE & Gate pass reps. & Robust OK \\
\midrule
encoded-state locally robust DGP & 25 & 10500 & 0.900 & 0.125 & 15 & 25 \\
\bottomrule
\end{tabular}
\end{table}

\begin{table}[H]
\centering\scriptsize
\caption{TD-CCP repeated-seed parameter uncertainty diagnostics.}
\begin{tabular}{lrrrrr}
\toprule
Parameter & Bias & RMSE & Emp. SD & Mean SE & 95\% cov. \\
\midrule
action\_1\_intercept & 0.008 & 0.080 & 0.081 & 0.069 & 0.960 \\
action\_1\_x0 & -0.056 & 0.095 & 0.078 & 0.089 & 0.920 \\
action\_1\_x1 & 0.045 & 0.118 & 0.112 & 0.076 & 0.880 \\
action\_2\_intercept & -0.058 & 0.122 & 0.109 & 0.092 & 0.920 \\
action\_2\_x0 & 0.026 & 0.136 & 0.136 & 0.116 & 0.920 \\
action\_2\_x1 & 0.076 & 0.145 & 0.125 & 0.108 & 0.800 \\
\bottomrule
\end{tabular}
\end{table}

\begin{table}[H]
\centering\scriptsize
\caption{TD-CCP repeated-seed inference diagnostics.}
\begin{tabular}{lrrrrr}
\toprule
Case & Cov. unit & Max zeta norm & Max lambda abs & Prelim. stat. & Prelim. pgrad \\
\midrule
encoded-state locally robust DGP & individual & 0.000 & 1.138 & 25 & 0.000 \\
\bottomrule
\end{tabular}
\end{table}

\subsection{Raw Neural Reward Diagnostic}
The raw neural flexible DGP is retained as a diagnostic failure record. It uses a frozen neural reward matrix with no finite true \(\theta\). TD-CCP is therefore not claimed to recover the raw reward matrix in this case; policy and counterfactual gates are reported separately from the finite-theta showcase.

Failed diagnostic gates: reward\_normalized\_rmse, value\_normalized\_rmse, q\_normalized\_rmse.

\begin{table}[H]
\centering\scriptsize
\caption{Raw neural reward diagnostic metrics.}
\begin{tabular}{lrrrrr}
\toprule
Case & Reward nRMSE & Policy TV & Value nRMSE & Q nRMSE & Gates fail \\
\midrule
raw-neural flexible diagnostic DGP & 0.801 & 0.009 & 0.113 & 0.120 & 3 \\
\bottomrule
\end{tabular}
\end{table}

\begin{table}[H]
\centering\scriptsize
\caption{Raw neural diagnostic counterfactual regret checks.}
\begin{tabular}{lrrr}
\toprule
Counterfactual & Policy TV & Value RMSE & Regret \\
\midrule
Type A & 0.009 & 0.004 & 0.004 \\
Type B & 0.013 & 0.012 & 0.012 \\
Type C & 0.007 & 0.003 & 0.003 \\
\bottomrule
\end{tabular}
\end{table}

\begin{table}[H]
\centering\scriptsize
\caption{Raw neural diagnostic gate audit.}
\begin{tabular}{lrrr}
\toprule
Gate & Value & Threshold & Pass \\
\midrule
converged & true & is true & true \\
reward\_normalized\_rmse & 0.801 & <= 0.100 & false \\
policy\_tv & 0.009 & <= 0.030 & true \\
value\_normalized\_rmse & 0.113 & <= 0.100 & false \\
q\_normalized\_rmse & 0.120 & <= 0.100 & false \\
type\_a\_regret & 0.004 & <= 0.050 & true \\
type\_b\_regret & 0.012 & <= 0.050 & true \\
type\_c\_regret & 0.003 & <= 0.050 & true \\
\bottomrule
\end{tabular}
\end{table}


The experiment simulates a firm value model where productivity $z \in [0,1]$ follows an AR(1) process with persistence $\rho = 0.8$.
Firms choose to operate ($a=1$) or exit ($a=0$) each period.
Flow utility is $u(1, z;\theta) = \theta_v z + \theta_{FC}$ and $u(0, z;\theta) = 0$, where $\theta_v = 5.0$ is the productivity value and $\theta_{FC} = -1.0$ is the fixed operating cost (negative sign).
The state space is represented by \tdNstates\ discrete bins covering $[0,1]$, which approximates the continuous process.
The dataset has \tdNobs\ observations from 3{,}000 individuals over 5 periods, with $\beta = \tdBeta$.

CCP must estimate the structural parameters by discretizing the state space into $K$ coarse bins and estimating the transition matrix $\hat P(z' \mid z, a)$ from data.
With $K = 5$ bins (width $0.2$), each bin averages over a wide range of true productivity values, introducing feature misspecification.
With $K = 20$ bins (width $0.05$), the features are more accurate but the transition matrix has 800 cells to estimate from the same data, producing noisy $\hat P$.
Neither discretization level escapes both sources of error simultaneously.
TD-CCP uses degree-8 polynomial basis functions over the normalized state $z_{\mathrm{norm}} \in [0,1]$ and never estimates the transition density.

Table~\ref{tab:tdccp_results} shows the results.
TD-CCP recovers $\theta_v$ within 1 percent and $\theta_{FC}$ within 7 percent of the true values, with valid robust standard errors reported in parentheses.
CCP with $K = 5$ bins achieves parameter RMSE of \ccpFiveRMSE, roughly 4 times worse than TD-CCP's \tdRMSE.
CCP with $K = 20$ bins achieves RMSE of \ccpTwentyRMSE, roughly 8 times worse, because the finer grid creates noisier transition estimates that compound the discretization error.
Behavioral cloning reaches \bcAcc\ policy accuracy but cannot recover structural parameters or produce standard errors.

The comparison illustrates the paper's central claim: there is no safe choice of bin width for CCP when the state is continuous.
Coarser bins introduce feature bias; finer bins introduce transition estimation variance.
TD-CCP sidesteps the dilemma by forming $h_k$ and $g$ directly from the observed next-state transitions $(z_t, a_t, z_{t+1})$, never forming $P(z' \mid z, a)$ at all.
The locally robust standard errors are valid at the parametric $\sqrt{n}$ rate despite the nonparametric polynomial first stage.

\begin{lstlisting}[basicstyle=\footnotesize\ttfamily]
from econirl.estimation.td_ccp import TDCCPEstimator, TDCCPConfig

config = TDCCPConfig(method="semigradient", basis_dim=8,
    ccp_method="logit", robust_se=True)
result = TDCCPEstimator(config=config).estimate(
    panel, utility, problem, transitions)
\end{lstlisting}

\paragraph{Code.}
\texttt{tdccp\_run.py} generates Table~\ref{tab:tdccp_results}. Estimator: \texttt{src/econirl/estimation/td\_ccp.py}.
\begin{verbatim}
.venv/bin/python papers/econirl_package/primers/tdccp/tdccp_run.py
\end{verbatim}

\end{document}
